% paper_25_hopf_gauge_structure.tex -- GeoVac Paper 25
% The Hopf Graph as a Lattice Gauge Structure: A Framework Observation
% Status: Draft, April 15, 2026
% Category: Synthesis (papers/synthesis/)
\documentclass[aps,pra,twocolumn,superscriptaddress,longbibliography]{revtex4-2}

\usepackage{amsmath,amssymb,amsthm,graphicx,xcolor,bm,braket,dcolumn,booktabs}
\usepackage{hyperref}

\newtheorem{theorem}{Theorem}
\newtheorem{observation}{Observation}
\newtheorem{conjecture}{Conjecture}
\newtheorem{proposition}{Proposition}

\begin{document}

\title{The Hopf Graph as a Lattice Gauge Structure:\\
A Framework Observation}
\thanks{Paper 25 in the Geometric Vacuum series}

\author{J.~Loutey}
\affiliation{Independent Researcher, Kent, Washington}

\date{April 15, 2026}

%% ========================================================================
%% ABSTRACT
%% ========================================================================

\begin{abstract}
The GeoVac framework places Coulomb bound-state quantum numbers on a
discrete graph that converges to the unit three-sphere $S^3$ in the
continuum limit.  This paper documents a framework observation: the
same graph, viewed through its incidence structure, instantiates the
combinatorial ingredients of a lattice gauge theory.  Nodes carry
scalar matter (0-forms), edges carry a $U(1)$ connection (1-forms),
and the Hopf quotient $S^3 \to S^2$ (collapsing the magnetic quantum
number $m_l$ along its $2l+1$ fiber) produces the discrete base graph
on which gauge-neutral propagation lives.  The node Laplacian
$L_0 = B B^{\top}$ is the matter propagator of the graph; the edge
Laplacian $L_1 = B^{\top} B$ is the discrete Hodge--1 Laplacian, whose
kernel counts independent cycles and whose structure is the
combinatorial analog of a photon propagator in temporal gauge.
Nothing about the mathematical decomposition is new: graph Hodge
theory (Lim; Schaub \emph{et al.}), lattice gauge theory (Wilson), and
the conformal invariance of 4D Maxwell (Cunningham; Bateman) are all
established separately.  What is new is the synthesis: the GeoVac
graph is a concrete case where all three structures coexist on a
single, physically motivated discretization.  Under the conjectural
interpretation of Paper~2, in which the bundle
$S^1 \to S^3 \to S^2$ is read as
(gauge phase)--(electron states)--(photon momenta), the
combinatorial formula $K = \pi(B + F - \Delta)$ becomes a
cross-sector spectral invariant of this lattice gauge structure.  We
do not claim a new derivation of the fine-structure constant.  We do
claim that the Paper~2 conjecture belongs inside an established
mathematical frame---one that makes its components
(selection principle, combination rule, circulant Hermiticity) into
gauge-theoretic statements rather than free-floating numerology.  We
close with open questions: whether the $S^5$ Bargmann--Segal lattice
of the 3D harmonic oscillator admits an analogous SU(3) gauge
structure, and whether the rank-2 Breit extension of Paper~22 lifts
the $U(1)$ 1-form picture to a higher-form gauge theory.
\end{abstract}

\maketitle

%% ========================================================================
%% I. INTRODUCTION
%% ========================================================================

\section{Introduction}
\label{sec:introduction}

The computational framework documented in Papers~0, 7, 14, 16, 18,
22, and 24 reformulates non-relativistic Coulomb bound-state physics
as spectral graph theory on a discretization of the unit
$S^3$~\cite{Paper0,Paper7,Paper14,Paper16,Paper18,Paper22,Paper24}.
The graph's nodes are quantum-number triples $(n, l, m_l)$, its edges
are the $L_\pm$ and $T_\pm$ ladder-operator transitions between
neighboring states, and its graph Laplacian converges to the
Laplace--Beltrami operator on the unit $S^3$ in the
continuum limit (Paper~7, eighteen symbolic proofs).

Separately, Paper~2 conjectures that the Hopf fibration
$S^1 \to S^3 \to S^2$---the same $S^3$ that hosts the electron's
bound-state manifold---is more than a geometric coincidence.  In the
Paper~2 reading, $S^3$ is the electron state space, $S^2$ is the
photon momentum sphere, and $S^1$ is the $U(1)$ gauge
phase~\cite{Paper2}.  The fine-structure constant is then
reinterpreted as the cost of projecting $S^3$ onto $S^2$, i.e.\ the
information absorbed in the fiber.  A combinatorial formula
\begin{equation}
K = \pi \left( B + F - \Delta \right)
\label{eq:K_intro}
\end{equation}
reproduces $\alpha^{-1}$ at $8.8 \times 10^{-8}$ with zero free
parameters and $p$-value $5.2 \times 10^{-9}$
(Paper~2~\cite{Paper2}).  Seven sprints of structural decomposition
(Phases 4B--4H) have identified each of $B$, $F$, $\Delta$ as
spectrally natural objects on distinct corners of the Hopf bundle,
but the additive combination rule itself remains unexplained.

This paper makes one observation.  The GeoVac Hopf graph is
simultaneously:
\begin{enumerate}
\item \emph{A discretization of $S^3$} (Fock projection; Paper~7).
\item \emph{A simplicial complex carrying a discrete Hodge
  decomposition} (Lim~\cite{Lim2020}; Schaub \emph{et
  al.}~\cite{Schaub2020}).
\item \emph{A lattice gauge structure} in the Wilson
  sense~\cite{Wilson1974}, with matter on nodes and a $U(1)$
  connection implicit on edges through the ladder-operator phases.
\end{enumerate}
No single paper in the literature has articulated (1)--(3) as one
object.  Lattice gauge theory has been formulated on flat hypercubic
lattices and, separately, on topologically nontrivial manifolds (S\({}^4\) for
instanton physics~\cite{BergLuscher1981,Luscher1982}, S\({}^3 \times
\mathbb{R}\) for finite-volume confinement
studies~\cite{Aharony2004}).  Graph Hodge theory has been applied to
signal processing and network analysis~\cite{Lim2020,Schaub2020}.
Conformal Maxwell theory on the conformal boundary of
AdS$_5$~\cite{EastwoodSinger1985,Maldacena1998} is a separate fourth
thread.  The communities have not overlapped on the specific case
where all four threads meet.

We do not claim a new derivation of $\alpha$.  We do claim that the
Paper~2 conjecture can be placed inside an established mathematical
framework.  In that framework, $B$ is a Casimir-weighted node-trace
on the $S^2$ quotient graph; $F$ is a spectral zeta of the infinite
Fock-degeneracy lattice at the packing exponent; $\Delta$ is a
boundary-mode count on the squared-Dirac sector.  The $\mathbb{Z}_3$
circulant structure of the cubic self-consistency equation
(Paper~2~\cite{Paper2}) becomes a three-fold projection structure on
the three bundle components.  What remains conjectural is not the
existence of these quantities---they are computable and have been
computed in Phases 4B--4G~\cite{Paper2}---but their additive
combination and its physical interpretation.

\paragraph*{Scope.} This is an observational synthesis paper.  No new
experimental or computational prediction is derived.  The purpose is
(i) to document a framework connection that has not been made in
either the GeoVac papers or the broader literature, (ii) to place the
conjectural components of Paper~2 inside an established mathematical
vocabulary, and (iii) to identify open questions (e.g., the $S^5$
Bargmann--Segal analog, the Breit tensor extension) that become
natural within the lattice-gauge framing.  The paper is a
companion to Paper~21's broader synthesis~\cite{Paper21}, focused on
one specific structural observation.

\paragraph*{What is conjectural and what is sharp.} The
\emph{mathematical} decomposition of the Hopf graph into
node-Laplacian, edge-Laplacian, and $S^2$-quotient components is
sharp: these are exactly computable, exactly described in Sections
II--III.  The \emph{physical} interpretation of the bundle components
(electron/photon/gauge) is Paper~2's conjecture, flagged throughout
as such.  The interpretive gap is deliberate: we do not assert the
conjecture is correct.  We assert that if it is correct, the
mathematics already in the literature places it in a specific
lattice-gauge context.

The remainder of this paper is organized as follows.  Section~II
surveys graph Hodge theory in the form we need (incidence matrix,
node/edge Laplacians, discrete Hodge decomposition, Betti numbers).
Section~III describes the GeoVac Hopf graph and its $S^2$ quotient in
this language.  Section~IV states the framework observation formally
and traces Paper~2's $(B, F, \Delta)$ ingredients to corresponding
spectral invariants.  Section~V discusses what the observation does
and does not predict.  Section~VI reviews related work and discusses
why the gap remained.  Section~VII lists open questions.  Section~VIII
concludes.

%% ========================================================================
%% II. GRAPH HODGE THEORY: THE TOOLS
%% ========================================================================

\section{Graph Hodge Theory: The Tools}
\label{sec:hodge}

This section collects the combinatorial machinery we use in
Sections~III--IV.  The content is standard; our contribution is the
identification in Section~III that the GeoVac graph is exactly the
kind of simplicial 1-complex to which this machinery applies.

\subsection{Incidence matrix and boundary operators}

Let $G = (V, E)$ be an undirected graph with $|V| = N_0$ nodes and
$|E| = N_1$ edges.  Fix an orientation on each edge (a choice, but
the Laplacians below are orientation-independent).  The signed
node-edge incidence matrix $B \in \mathbb{R}^{N_0 \times N_1}$ has
entries
\begin{equation}
B_{iv,\, e} =
\begin{cases}
+1 & \text{if edge $e$ starts at node $v$},\\
-1 & \text{if edge $e$ ends at node $v$},\\
0  & \text{otherwise}.
\end{cases}
\label{eq:incidence}
\end{equation}
Geometrically, $B$ is the discrete exterior derivative from 0-cochains
(scalar functions on nodes) to 1-cochains (scalar functions on
oriented edges), $d_0 = B^{\top}$, and its transpose $B$ is the
discrete boundary operator $\partial_1$ from 1-chains to 0-chains.

\subsection{Node and edge Laplacians}

The combinatorial Hodge--0 and Hodge--1 Laplacians
are~\cite{Lim2020,Schaub2020}
\begin{align}
L_0 &= B B^{\top} \in \mathbb{R}^{N_0 \times N_0},
\label{eq:L0}\\
L_1 &= B^{\top} B \in \mathbb{R}^{N_1 \times N_1}.
\label{eq:L1}
\end{align}
$L_0$ is the familiar node Laplacian $D - A$
(degree matrix minus adjacency matrix).  $L_1$ is the edge Laplacian,
acting on functions defined on edges.  If the simplicial complex has
triangular 2-faces with incidence matrix $B_2$, the full Hodge--1
Laplacian acquires an upper piece $B_2 B_2^{\top}$; we work with
triangle-free simplicial 1-complexes throughout this paper, so
$L_1 = B^{\top} B$ is the complete edge Laplacian.

\subsection{Discrete Hodge decomposition}

The cochain space $\mathbb{R}^{N_1}$ (edge functions) decomposes
orthogonally:
\begin{equation}
\mathbb{R}^{N_1} = \operatorname{im}(d_0) \oplus \ker(L_1).
\label{eq:hodge_decomp}
\end{equation}
The first summand contains \emph{gradient} edge functions,
$\omega = d_0 f = B^{\top} f$ for some node function $f$---the edge
function is ``pure gauge'' in the sense that it is a gradient.  The
second summand, $\ker(L_1)$, contains the \emph{harmonic} 1-cochains:
edge functions that are not gradients but whose Hodge--1 Laplacian
vanishes.  By the combinatorial Hodge theorem, the dimension of the
harmonic subspace equals the first Betti number of the graph:
\begin{equation}
\dim \ker(L_1) = \beta_1 = N_1 - N_0 + c,
\label{eq:betti1}
\end{equation}
where $c$ is the number of connected components.  $\beta_1$ counts
independent cycles: edge functions that circulate without being
gradients.  These are the combinatorial analog of a non-trivial
holonomy along a loop, i.e., the analog of a Wilson loop in the
Wilson formulation of lattice gauge theory.

\subsection{Node versus edge spectra}

The nonzero eigenvalues of $L_0 = B B^{\top}$ and $L_1 = B^{\top} B$
coincide by the singular-value decomposition: if $B = U \Sigma V^{\top}$
then $L_0 = U \Sigma^2 U^{\top}$ and $L_1 = V \Sigma^2 V^{\top}$.
The edge Laplacian therefore contributes no new eigenvalues, only new
kernel dimensions (the $\beta_1$ harmonic 1-forms beyond the
$c$ harmonic 0-forms).  This relation, elementary in matrix algebra,
will matter in Section~IV: the edge Laplacian's \emph{topological}
content (counting cycles) is distinct from its \emph{spectral}
content (shared with the node Laplacian).

\subsection{Lattice gauge theory dictionary}

The ingredients of Wilson's~\cite{Wilson1974} lattice gauge theory
map directly to the graph-Hodge ingredients:
\begin{center}
\begin{tabular}{lll}
\toprule
Wilson & Graph Hodge & Role \\
\midrule
Site $x$            & Node $v$         & Matter location \\
Link $(x, x+\hat{\mu})$ & Edge $e$     & Gauge-field location \\
Link variable $U_e \in G$ & Edge phase $\theta_e \in [0, 2\pi)$ & $U(1)$ connection \\
Plaquette $P$       & Triangle/cycle   & Curvature / $F_{\mu\nu}$ \\
Wilson loop $W_C$   & Holonomy $\prod_{e \in C} e^{i\theta_e}$ & Gauge-invariant observable \\
Plaquette action $\mathrm{Re}\,\mathrm{tr}\,U_P$ & Curvature norm $|F|^2$ & Gauge kinetic term \\
Gauge transformation & $\theta_e \to \theta_e + (\phi_{t(e)} - \phi_{h(e)})$ & 0-coboundary shift \\
Matter propagator  & $(L_0 + m^2)^{-1}$ & Covariant Laplacian \\
Photon propagator  & $(L_1 + m^2)^{-1}$ & Hodge--1 resolvent \\
\bottomrule
\end{tabular}
\end{center}

The match is exact at the level of combinatorial data.  The only
question is whether a given graph carries the physical content that
makes this dictionary nontrivial.  For a random graph with arbitrary
$U(1)$ phases, the dictionary is empty: there is no physical
interpretation of the gauge degrees of freedom.  For the GeoVac Hopf
graph, the ladder-operator structure provides a canonical assignment
of edge phases through the angular-momentum and radial transition
amplitudes.  Section~III makes this explicit.

%% ========================================================================
%% III. THE GEOVAC HOPF GRAPH
%% ========================================================================

\section{The GeoVac Hopf Graph}
\label{sec:hopfgraph}

\subsection{Nodes and edges}

The GeoVac graph at cutoff $n_{\max}$ has nodes labeled by
$(n, l, m_l)$ with $n = 1, \ldots, n_{\max}$, $l = 0, \ldots, n-1$,
and $m_l = -l, \ldots, +l$.  The total node count is
$N_0(n_{\max}) = \sum_{n=1}^{n_{\max}} n^2$.  Edges come in two
kinds~\cite{Paper7,Paper14}:
\begin{itemize}
\item \textbf{Angular} edges connect
  $(n, l, m_l) \leftrightarrow (n, l, m_l \pm 1)$ at fixed
  $(n, l)$---these are the $L_\pm$ transitions of the orbital
  angular-momentum ladder.
\item \textbf{Radial} edges connect
  $(n, l, m_l) \leftrightarrow (n', l, m_l)$ where $n' = n \pm 1$
  (and sometimes, depending on the discretization scheme, with
  $\Delta l = \pm 1$)---these are the $T_\pm$ hydrogenic radial
  transitions.
\end{itemize}
At $n_{\max} = 3$ the graph has $N_0 = 1 + 4 + 9 = 14$ nodes and
$N_1 = 13$ edges, with $c = 3$ connected components (s-chain,
p-block, d-chain).  The first Betti number is
$\beta_1 = N_1 - N_0 + c = 13 - 14 + 3 = 2$~\cite{Q1B}.  The two
independent cycles both live in the p-block, where the angular
$m_l$-ladder and the radial $n$-transition together form square-like
loops.

\subsection{The Hopf quotient: $S^3 \to S^2$}

The Hopf fibration $S^1 \to S^3 \to S^2$ projects the $S^3$ onto the
base $S^2 \cong \mathbb{CP}^1$, collapsing each $S^1$ fiber to a
point.  On the GeoVac graph, the natural discrete implementation of
this quotient collapses each magnetic fiber $\{(n, l, m_l) :
m_l = -l, \ldots, +l\}$ at fixed $(n, l)$ to a single sector label
$(n, l)$.  The Hopf base is then a graph on $(n, l)$ sectors.

At $n_{\max} = 3$, the base graph has 6 sectors,
\begin{equation}
\{(1,0),\, (2,0),\, (2,1),\, (3,0),\, (3,1),\, (3,2)\},
\label{eq:s2_sectors}
\end{equation}
with sector weights $(2l+1) = 1, 1, 3, 1, 3, 5$ (the fiber
dimensions).  Inter-sector edges of the base graph count the number
of $S^3$ edges crossing between sectors; intra-sector angular edges
collapse into ``fiber'' structure that does not appear in the base.
Track~$\alpha$-D~\cite{Paper2} computed this base graph explicitly:
its Laplacian has eigenvalues $\{0, 0, 0, 1, 3, 6\}$ (three zero
modes reflecting the three disconnected components inherited from
the $S^3$ graph).

\subsection{$U(1)$ phase on edges}

The $L_\pm$ and $T_\pm$ ladder operators act on the spherical
harmonic basis with computable matrix elements.  For angular edges,
\begin{equation}
\braket{n, l, m_l + 1 | L_+ | n, l, m_l}
= \hbar \sqrt{l(l+1) - m_l(m_l+1)} \, e^{i\phi_{L_+}}.
\label{eq:Lplus}
\end{equation}
For radial edges, the amplitudes are the hydrogenic Clebsch--Gordan
coefficients times a phase.  Under a global $U(1)$ rotation of the
wavefunctions $\psi_{n,l,m_l} \to e^{i\chi}\psi_{n,l,m_l}$, the edge
amplitudes pick up no phase; under a node-dependent
$\psi \to e^{i\chi_v}\psi$, the edge amplitude for $e = (v, v')$ picks
up the phase $e^{i(\chi_{v'} - \chi_v)}$---this is exactly the
$U(1)$ gauge transformation in the Wilson dictionary of Section~II.E.

The GeoVac graph is therefore naturally equipped with a $U(1)$ bundle:
each edge carries a complex phase, and the phases transform
covariantly under node-local $U(1)$ redefinitions of the wavefunctions.
The curvature of this bundle is the holonomy around each closed
cycle.  At $n_{\max} = 3$, $\beta_1 = 2$ cycles (both in the
p-block) carry nontrivial holonomy.

\subsection{Node and edge Laplacians of the GeoVac graph}

With the incidence matrix $B$ as in Eq.~(\ref{eq:incidence}), the
node Laplacian $L_0 = B B^{\top}$ is the one whose continuum limit
is $-\Delta_{S^3}$ (Paper~7).  The edge Laplacian
$L_1 = B^{\top} B$ has explicit spectrum at $n_{\max} = 3$
computed in Track Q1--B~\cite{Q1B}:
\begin{equation}
\mathrm{spec}(L_1) = \{0^{(2)},\, \tfrac{3-\sqrt{5}}{2},\, 1^{(2)},\,
  \tfrac{5-\sqrt{5}}{2},\, 2,\, \tfrac{3+\sqrt{5}}{2},\, 3^{(3)},\,
  \tfrac{5+\sqrt{5}}{2},\, 5\}.
\label{eq:L1_spectrum}
\end{equation}
The two zero modes confirm $\beta_1 = 2$.  The golden-ratio
eigenvalues come from the path graph $P_5$ of the d-chain; the
eigenvalues $\{1, 3, 3, 5\}$ come from the p-block (which contains
both cycles).  The characteristic polynomial factors completely over
$\mathbb{Q}(\sqrt{5})$: all edge-Laplacian spectral invariants are
algebraic numbers with denominator dividing $2^k$ times a power of
$\sqrt{5}$.

Observation~1 follows directly:

\begin{observation}[Edge Laplacian is algebraic on a finite
    Hopf graph]
\label{obs:edge_algebraic}
The edge Laplacian $L_1 = B^{\top} B$ of the GeoVac graph at any
finite $n_{\max}$ has integer-valued entries, hence algebraic
eigenvalues.  No transcendental spectral invariant can live on the
edge Laplacian of a finite Hopf graph.  Transcendentals enter only
through spectral zeta functions in the continuum limit
($n_{\max} \to \infty$), consistent with the Phase~4C--4D finding
that $F = \pi^2/6$ is not accessible from any finite-graph
Hopf-fiber trace~\cite{Paper2}.
\end{observation}

\subsection{The matter and gauge propagators}

Under the Paper~2 interpretation:
\begin{itemize}
\item $L_0 = B B^{\top}$ is the matter propagator's inverse: bound
  electron states propagate by the $S^3$ Laplace--Beltrami operator
  in the continuum.
\item $L_1 = B^{\top} B$ is the gauge-sector propagator's inverse:
  the $U(1)$ photon-like edge modes, in temporal gauge, propagate by
  the discrete Hodge--1 Laplacian.  The kernel dimension
  $\beta_1$ counts independent Wilson-loop classes.
\end{itemize}
This is the observational claim of the paper.  We now formalize it.

%% ========================================================================
%% IV. THE FRAMEWORK OBSERVATION
%% ========================================================================

\section{The Framework Observation, Formalized}
\label{sec:framework}

\begin{proposition}[Lattice gauge structure of the GeoVac graph]
\label{prop:gauge_structure}
Let $G$ be the GeoVac graph at cutoff $n_{\max}$, with nodes labeled
by $(n, l, m_l)$, edges labeled by angular and radial ladder-operator
transitions, and signed incidence matrix $B$.  Then:
\begin{enumerate}
\item $G$ is a triangle-free simplicial 1-complex, so its
  combinatorial Hodge Laplacians are $L_0 = B B^{\top}$ and
  $L_1 = B^{\top} B$.
\item The node Laplacian $L_0$ converges to the Laplace--Beltrami
  operator $-\Delta_{S^3}$ in the continuum limit
  (Paper~7~\cite{Paper7}).
\item The edge Laplacian $L_1$ has kernel dimension
  $\beta_1(G) = N_1 - N_0 + c$; these harmonic 1-forms are the
  discrete analog of topologically nontrivial $U(1)$ connections.
\item The Hopf quotient (collapsing $m_l$-fibers) yields a
  sub-graph $G^{S^2}$ on $(n, l)$ sectors, with Laplacian spectrum
  $\{0, 0, 0, 1, 3, 6\}$ at $n_{\max} = 3$ (Track
  $\alpha$-D~\cite{Paper2}).
\item Natural $U(1)$ phases on edges arise from the
  ladder-operator amplitudes; a node-local redefinition
  $\psi_v \to e^{i\chi_v} \psi_v$ acts as a gauge transformation in
  the Wilson sense on these phases.
\end{enumerate}
\end{proposition}

This proposition collects facts already scattered across
Papers~0, 2, 7, 14, and the Phase~4B Track $\alpha$-D data set.  The
new content is the unified \emph{reading} of those facts.  The reading
is not a new computational claim; it is a dictionary that places the
GeoVac graph in the Wilson--Hodge frame.

\subsection{Where Paper~2's $(B, F, \Delta)$ live in this frame}

The three structural ingredients of Paper~2's conjectural combination
rule $K = \pi(B + F - \Delta)$ have, by the Phase~4B--4H sprint
program, been localized on distinct components of the Hopf bundle.
In the present framework, the localizations acquire gauge-theoretic
labels.

\paragraph*{$B = 42$: Casimir-weighted node-trace on the $S^2$
quotient.}  The Phase~4B track $\alpha$-D finding~\cite{Paper2} is
that
\begin{equation}
B(n_{\max}) = \sum_{n=1}^{n_{\max}} \sum_{l=0}^{n-1}
  (2l+1)\,l(l+1)
\label{eq:B_casimir}
\end{equation}
is the trace of the orbital-angular-momentum Casimir $\hat{L}^2$,
weighted by the fiber dimension $2l+1$, summed over the sectors of
$G^{S^2}$.  In lattice-gauge language, $B$ is a \emph{matter
operator expectation value on the $S^2$ base}: the photon propagates
on the sectors (the $(n, l)$ labels), and the Casimir
$\hat{L}^2$ is a gauge-invariant operator diagonal in this basis.
$B = 42$ at $n_{\max} = 3$ is the value of this operator summed over
the base nodes, weighted by their $U(1)$ fiber dimensions.

\paragraph*{$F = \pi^2 / 6$: Dirichlet zeta of the infinite
Fock-degeneracy lattice.}  The Phase~4F Track $\alpha$-J
result~\cite{Paper2} is that
\begin{equation}
F = \sum_{n=1}^{\infty} \frac{n^2}{n^4}
  = \zeta_R(2) = \frac{\pi^2}{6}.
\label{eq:F_zeta}
\end{equation}
The weight $g_n = n^2$ is the Fock degeneracy of the $n$-th $S^3$
shell; the exponent $s = 4$ is the packing exponent
$d_{\max} = \dim(\mathbb{R}^4)$ from Paper~0.  This is not a
finite-graph spectrum: it is the analytic continuation of a graph
spectral sum to the full infinite Hopf lattice.  In gauge-theoretic
language, $F$ is a \emph{regularized gauge-field partition function}
on the infinite base: a spectral zeta of the photon-sector degrees
of freedom summed over all shells.

\paragraph*{$\Delta = 1/40$: Dirac-mode boundary count.} The Phase~4H
Track SM-D result~\cite{Paper2} is that
\begin{equation}
\Delta^{-1} = g_3^{\mathrm{Dirac}}(S^3)
  = 2(n+1)(n+2)\bigg|_{n=3} = 40.
\label{eq:Delta_dirac}
\end{equation}
This is the single-chirality Dirac mode count of the third
Camporesi--Higuchi eigenspace on the unit $S^3$.  In
lattice-gauge-coupled-to-matter language, $\Delta$ is a
\emph{boundary-mode correction} counting the charged-spinor states at
the truncation edge.  Its entry with a minus sign in the combination
$K = \pi(B + F - \Delta)$ is consistent (under the Paper~2 reading)
with a vacuum-polarization correction that reduces, rather than
enhances, the effective coupling at the cutoff.

\subsection{What the framework observation claims}

\begin{observation}[Framework observation]
\label{obs:framework}
Under the Paper~2 conjecture identifying $S^1 \to S^3 \to S^2$ with
(gauge phase)--(electron states)--(photon momenta), the Paper~2
combination rule $K = \pi(B + F - \Delta)$ is an additive combination
of (i)~a finite matter-operator trace on the $S^2$ base, (ii)~an
infinite gauge-field spectral zeta on the Fock lattice, and (iii)~a
boundary-mode count on the charged-spinor spectrum.  All three
quantities are intrinsic spectral invariants of the discrete
lattice-gauge structure of the Hopf graph.
\end{observation}

\begin{observation}[What the framework observation does \emph{not}
claim]
\label{obs:not_claim}
The framework observation does not derive the additive combination
rule itself.  It does not explain why the three invariants combine
with a common prefactor $\pi$, nor why $F$ enters with a plus sign
and $\Delta$ with a minus.  The combination rule remains an
empirical numerical match at $8.8 \times 10^{-8}$, with $p$-value
$5.2 \times 10^{-9}$ against an exhaustive combinatorial search
(Paper~2)~\cite{Paper2}, and its status as conjecture is unchanged.
\end{observation}

\subsection{Circulant structure reread}

Paper~2 shows that the cubic self-consistency equation
$\alpha^3 - K\alpha + 1 = 0$ is the characteristic polynomial of a
traceless $\mathbb{Z}_3$-symmetric circulant
matrix~\cite{Paper2}.  The $\mathbb{Z}_3$ action
permutes three abstract coordinates; in Paper~2's interpretation,
these are $(B, F, \Delta)$ or the three Hopf-bundle components.  The
$U(1)$ gauge structure forces Hermiticity of the circulant, which in
turn forces the specific product relation $bc = K/3$ between
off-diagonal entries.

In the present framework, the $\mathbb{Z}_3$ symmetry is the cyclic
permutation of the three distinct \emph{tiers} identified in
Section~IV.A: a finite matter trace, an infinite gauge zeta, and a
boundary Dirac mode count.  The Hermiticity arises from the
compatibility of the $U(1)$ gauge structure across the three tiers.
This is the same observation as Paper~2, rewritten in
lattice-gauge language.  It does not sharpen the conjecture; it does
label its components with established gauge-theoretic names.

%% ========================================================================
%% V. WHAT THIS PREDICTS AND DOES NOT PREDICT
%% ========================================================================

\section{What This Predicts, and What It Does Not}
\label{sec:predictions}

\subsection{What this does \emph{not} produce}

This paper does \emph{not}:
\begin{itemize}
\item Derive the fine-structure constant $\alpha$ from first
  principles.  Paper~2 is, and remains, a conjecture.  The present
  paper shifts the conjecture into gauge-theoretic language without
  removing any of its undetermined inputs.
\item Prove the combination rule $K = \pi(B + F - \Delta)$.  Phases
  4B--4H eliminated nine specific derivation mechanisms; the
  additive rule remains a numerical coincidence at $p = 5.2 \times
  10^{-9}$~\cite{Paper2}.
\item Produce a continuum limit in which the edge Laplacian
  $L_1$ of the GeoVac graph converges to the Hodge--1 Laplacian on
  the unit $S^3$.  Such a convergence is a mathematically
  well-posed but computationally nontrivial question; at present, we
  have only the finite-$n_{\max}$ edge-Laplacian spectrum
  (Section~III.D).
\item Give the photon dynamical content.  The edge Laplacian is a
  combinatorial object; no dispersion relation, Lorentz structure,
  or QED vertex emerges from it alone.
\end{itemize}

\subsection{What the framework observation does make available}

What the observation does provide is a set of \emph{natural next
computations} that become well posed in the lattice-gauge frame but
would have been ad hoc before:

\paragraph*{Continuum-limit edge-Laplacian convergence.} If the
GeoVac graph node Laplacian $L_0$ converges to $-\Delta_{S^3}$ in
the continuum (Paper~7), does $L_1$ converge to the
Hodge--1 Laplacian on $S^3$?  The continuous Hodge--1 spectrum on
$S^3$ is known exactly: eigenvalues $n(n+2)$ with degeneracy
$2n(n+2)$, $n \geq 1$, and zero $\beta_1$
(Ikeda--Taniguchi~\cite{IkedaTaniguchi1978}).  The finite-$n_{\max}$
graph has $\beta_1 = 2$ at $n_{\max} = 3$, $\beta_1 = ?$ at higher
$n_{\max}$ (monotone or asymptotic behavior computable).  If
$\beta_1 / N_1 \to 0$ and the nonzero spectrum converges to
$n(n+2)$, we have a discretization of the Hodge--1 structure on
$S^3$ compatible with the Hodge--0 convergence of Paper~7.

\paragraph*{Higher Hopf bundles.} The $S^5$ Bargmann--Segal lattice
of Paper~24~\cite{Paper24} is built on the Hardy space
$\mathcal{H}^2(S^5)$ of the holomorphic SU(3) $(N, 0)$
representations.  $S^5 \to \mathbb{CP}^2$ is the next complex Hopf
fibration after $S^3 \to S^2 = \mathbb{CP}^1$.  Does the
Bargmann--Segal graph carry an analogous SU(3) non-abelian gauge
structure?  Candidate check: the SU(3) Casimir values on the
$(N, 0)$ irreps, weighted by fiber dimensions and summed under an
analogous selection principle, produce a non-abelian analog of
$B = 42$.  Paper~24's HO rigidity theorem~\cite{Paper24} already
rules out calibration-$\pi$ content on the $S^5$ graph, consistent
with the expectation that a non-abelian SU(3) gauge-like structure
would give a structurally different (likely pure-rational)
combination rule, rather than an $\alpha$-like transcendental.

\paragraph*{Rank-2 Breit extension.} Paper~22's Sprint 2 extension
to the rank-2 Breit interaction opens a second family of edges
carrying rank-2 tensor content~\cite{Paper22}.  In the Hodge
vocabulary, this is a 2-cochain structure built on top of the
1-cochain $U(1)$ connection.  A natural question: does the combined
1-form + 2-form simplicial structure carry a higher-form gauge
theory (e.g., a 2-form $B$-field of the kind familiar in string
compactifications)?  This is speculative; we flag it only to point
out that the lattice-gauge frame opens a coherent question where
previously there was just an observation about integral sparsity.

\paragraph*{Paper~2 tests independent of $\alpha$.} If the Paper~2
interpretation holds, the $U(1)$-Hopf structure of the GeoVac graph
should show up in observables beyond $\alpha$.  Candidates within
the GeoVac pipeline: (i)~systematic relations between the spin-orbit
splittings computed in Tier~2 (Paper~14 Sec.~V~\cite{Paper14}) and
the edge-Laplacian structure of the spinor-bundle graph of
Paper~18 Sec.~IV~\cite{Paper18}; (ii)~the rank-2 Breit tensor
contribution to ERI density (Paper~22 Sprint 2
extension~\cite{Paper22}) interpreted as a 2-form contribution to
the gauge-sector Hamiltonian.  These are testable against the
existing GeoVac codebase and give $\alpha$-independent pressure on
the Paper~2 interpretation.  We leave the tests to future sprints.

%% ========================================================================
%% VI. RELATED WORK
%% ========================================================================

\section{Related Work}
\label{sec:relatedwork}

\subsection{Lattice gauge theory on curved manifolds}

Wilson~\cite{Wilson1974} formulated lattice gauge theory on flat
hypercubic lattices in Euclidean $\mathbb{R}^4$, with matter on sites
and gauge fields on links.  Luscher~\cite{Luscher1982} developed the
topology of lattice gauge fields on discretized $S^4$, with
continuity constraints ensuring a well-defined integer topological
charge; Berg and Luscher~\cite{BergLuscher1981} used $S^4$ lattices
for instanton studies.  These are \emph{compactification} strategies,
not claims that $S^3$ or $S^4$ is a \emph{natural} geometry for the
physics.  Aharony \emph{et al.}~\cite{Aharony2004} formulated large-$N$
Yang--Mills on $S^3 \times \mathbb{R}$ as a probe of the
Hagedorn/deconfinement phase transition, again as a
finite-volume/infrared-regulator choice rather than as a natural
arena.  The category distinction matters: these works treat the
manifold as a regulator; the GeoVac framework treats it as the
physical arena of the electron bound states (via Fock's
projection~\cite{Fock1935}), \emph{and} as a geometric regulator of
the gauge degrees of freedom.

Witten~\cite{Witten1989} formulated Chern--Simons theory on $S^3$ and
computed the Jones polynomial of knots from its partition function.
This is a \emph{topological} gauge theory in three spacetime
dimensions without local dynamical degrees of freedom.  The GeoVac
lattice-gauge reading would be a dynamical 3+1-dimensional theory
(matter + photon propagation); the shared vocabulary of
``connections on $S^3$'' is suggestive but not structurally identical.

Recent work (post-2015) on simplicial lattice gauge theory
(e.g., work on tensor-network formulations and categorical
symmetries) has developed abstract combinatorial gauge theories on
simplicial complexes, but without the specific physical motivation
(Fock projection) that the GeoVac graph carries.

\subsection{Graph Hodge theory}

Lim~\cite{Lim2020} and Schaub \emph{et al.}~\cite{Schaub2020}
articulated the combinatorial Hodge theory of graphs and simplicial
complexes in modern form, motivated by applications in signal
processing, ranking, and random walks.  The mathematical framework is
standard (going back to Eckmann in the 1940s and to algebraic
topology generally), but its packaging in a form usable by
engineers and data scientists is recent.  The Hodge Laplacian
$L_1 = B^{\top} B + B_2 B_2^{\top}$, its kernel as $\beta_1$, and
the exact/coexact/harmonic decomposition of edge functions are all
standard content of these surveys.  The gap is that these surveys do
not connect to lattice gauge theory or to quantum field theory on
graphs; the communities are disjoint.

\subsection{Conformal Maxwell theory and the AdS/CFT boundary}

The conformal invariance of source-free Maxwell theory in 4D is a
classical result of Cunningham~\cite{Cunningham1909} and
Bateman~\cite{Bateman1910}.  Under conformal compactification of
Minkowski space, the conformal boundary is $S^3 \times \mathbb{R}$,
which is also the conformal boundary of AdS$_5$
(Maldacena~\cite{Maldacena1998}).  Eastwood and
Singer~\cite{EastwoodSinger1985} developed a conformally invariant
Maxwell gauge on this boundary; the photon modes expand in vector
spherical harmonics on $S^3$.  The continuous Hodge--1 spectrum on
$S^3$ (Ikeda--Taniguchi~\cite{IkedaTaniguchi1978}) is the relevant
continuum limit.

Under the Paper~2 interpretation, the GeoVac $S^3$ (from Fock's
momentum-space projection of the electron) and the
conformal-compactification $S^3$ (from the conformal boundary of
AdS$_5$) are \emph{distinct} manifolds that happen to carry the same
topology.  Whether they can be unified into a single mathematical
object---so that one $S^3$ carries both matter and gauge content in
a coherent way---is an open foundational question; it is not
resolved by the present observation.  The present observation is
more modest: the GeoVac graph, as a finite combinatorial object,
carries both matter and gauge content simultaneously via its
node-Laplacian / edge-Laplacian structure, regardless of how the
continuum limit is ultimately interpreted.

\subsection{Why the gap remained}

Three communities would have had to overlap for the
framework observation to be routine:
\begin{enumerate}
\item \emph{Lattice gauge theorists}, who know Wilson's formulation
  but work on flat hypercubic lattices, not on Fock-projected
  momentum-space $S^3$ lattices.
\item \emph{Graph Hodge theorists}, who know $L_0$/$L_1$ but work
  on engineering applications (ranking, signal processing), not on
  physical lattice gauge theory.
\item \emph{Atomic/quantum chemistry theorists}, who know Fock's
  1935 projection but treat it as a solved analytical curiosity, not
  as a computational framework underlying a discrete gauge theory.
\end{enumerate}
The GeoVac project arrived at (3) for purely computational reasons
(sparse qubit Hamiltonians from angular-momentum selection rules;
see Paper~14~\cite{Paper14}).  The overlap with (1) and (2) was
invisible from inside any single community.  Recognizing it from
within the GeoVac pipeline required the Paper~2 conjecture, which in
turn required the Phase~4B--4H structural decomposition to
identify the three spectral homes of $B$, $F$, $\Delta$.  We flag
the gap to place future GeoVac work at the intersection
deliberately, rather than leaving the connection implicit.

%% ========================================================================
%% VII. OPEN QUESTIONS
%% ========================================================================

\section{Open Questions}
\label{sec:open}

\subsection{SU(3) gauge structure on the Bargmann--Segal $S^5$ lattice}

Paper~24~\cite{Paper24} constructs the Bargmann--Segal lattice
of the 3D isotropic harmonic oscillator on the Hardy space
$\mathcal{H}^2(S^5)$ of SU(3) $(N, 0)$ symmetric representations.
Its nodes are labeled $(N, l, m_l)$, its edges are the
$\Delta N = \pm 1$, $\Delta l = \pm 1$ dipole transitions.  The
underlying bundle is $S^5 \to \mathbb{CP}^2$, the next complex Hopf
fibration.

If the present framework observation generalizes, the Bargmann--Segal
graph carries an SU(3) non-abelian gauge structure: nodes carry
matter, edges carry SU(3)-valued connections, and the $\mathbb{CP}^2$
quotient is the base of the non-abelian bundle.  A natural first
check is whether an analog of the $B = 42$ Casimir-weighted trace
exists with the SU(3) quadratic Casimir replacing the SO(3) orbital
Casimir, and whether an analog of the selection principle
$B/N = \dim = 5$ holds at any natural cutoff.  Paper~24's HO rigidity
theorem~\cite{Paper24} already establishes that the Bargmann--Segal
graph is bit-exactly $\pi$-free in exact rational arithmetic at every
$N_{\max}$; this is consistent with an SU(3) gauge-like structure
producing only rational invariants at finite $N_{\max}$, with
transcendental content (if any) entering only in the infinite limit,
as it does for the $S^3$ Hopf case via the infinite Dirichlet series
of Paper~4F~\cite{Paper2}.

Whether the non-abelian case produces a physical coupling constant
(e.g., QCD $\alpha_s$ at some reference scale) is a highly speculative
question that we flag but do not investigate.  The Phase~4H negative
result~\cite{Paper2} showed that the naive generalization of
Paper~2's selection principle to higher Hopf bundles does not produce
the Standard Model couplings.  Any non-abelian extension must avoid
reproducing that known-negative path.

\subsection{Rank-2 Breit tensor: a 2-form gauge structure?}

Paper~22's Sprint 2 extension treats the rank-2 Breit interaction as
a genuinely new tensor sector of the GeoVac graph, with sparsity
density depending on $l_{\max}$ but structurally distinct from the
rank-0 Coulomb sector~\cite{Paper22}.  In the present framework, the
rank-0 Coulomb edges are 1-cochains (a $U(1)$ connection on nodes);
the rank-2 Breit edges would be 2-cochains (a 2-form on faces or
pairs of edges).

A 2-cochain sector on the GeoVac graph would correspond (in the
continuum) to a $B$-field of the kind familiar from string
compactifications, or to the 2-form part of a higher-form gauge
theory.  Its combinatorial Laplacian is $L_2 = B_2^{\top} B_2 +
B_3 B_3^{\top}$, if triangle and tetrahedral faces are included.
Whether the Breit tensor structure closes into a well-defined
higher-form gauge theory on the GeoVac graph---or whether it is
better viewed as a nonlinear coupling on the $U(1)$ sector---is an
open question.  The Sprint 2 computational data (ERI densities at
rank-2 $l_{\max}$ values) are the input to such a test; the test
itself has not been performed.

\subsection{Observables beyond $\alpha$}

The Paper~2 interpretation, if correct, should have implications
beyond the fine-structure constant.  Candidate observables to test:
\begin{itemize}
\item \emph{Magnetic moment anomalies.}  The $U(1)$ Hopf structure
  places the $g$-factor in a specific combinatorial context (via
  the spinor-bundle analog of the $U(1)$ matter-gauge coupling).
  Paper~14 Tier~2~\cite{Paper14} computes relativistic spin-orbit
  splittings; whether these decompose into Hopf-base and
  Hopf-fiber contributions with the same $\mathbb{Z}_3$ circulant
  structure as Paper~2 is computable.
\item \emph{$\alpha$-dependence at finite $n_{\max}$.}  If
  $K = \pi(B + F - \Delta)$ is indeed the lattice-gauge invariant
  producing $\alpha^{-1}$, then at $n_{\max} \neq 3$ the truncated
  $B(n_{\max})$ and $\Delta(n_{\max})$ differ and give different
  ``$\alpha$'' values.  The selection principle $B/N = 3$
  uniquely at $n_{\max} = 3$ says that only this cutoff reproduces
  the physical $\alpha$.  This is a prediction: no other
  $n_{\max}$ produces a recognizable physical coupling.  Paper~2's
  negative result on higher Hopf bundles is consistent with this
  selection criterion; a systematic $n_{\max}$ scan is the obvious
  next test.
\item \emph{Sharpness of the combination rule.}  Under what class of
  modifications to $G$ (e.g., edge reweightings, adding or removing
  specific transitions) does $K$ remain within $10^{-7}$ of
  $\alpha^{-1}$?  A combinatorial robustness check would distinguish
  a genuine spectral invariant (stable under small perturbations)
  from a fine-tuned numerical coincidence (collapsing under any
  perturbation).
\end{itemize}
None of these tests has been performed.  The present paper's
contribution is to identify them as well-posed within the framework,
not to execute them.

%% ========================================================================
%% VIII. CONCLUSION
%% ========================================================================

\section{Conclusion}
\label{sec:conclusion}

The GeoVac framework builds a discrete graph on the unit $S^3$
whose node Laplacian converges to the Laplace--Beltrami operator on
$S^3$ in the continuum limit.  Section~III of this paper made the
observation---previously implicit in the Paper~2 conjecture but never
stated as its own claim---that the same graph, together with its
signed incidence matrix $B$, carries all the combinatorial
ingredients of a lattice gauge theory.  The node Laplacian
$L_0 = B B^{\top}$ is the matter propagator; the edge Laplacian
$L_1 = B^{\top} B$ is the discrete Hodge--1 Laplacian; the Hopf
quotient collapsing $m_l$-fibers produces the $S^2$ base graph on
which gauge-neutral propagation lives.

Under the Paper~2 conjectural interpretation, the three components
of the Hopf bundle acquire physical labels: $S^3$ as electron bound
states, $S^1$ as $U(1)$ gauge phase, $S^2$ as photon momentum sphere.
The Paper~2 combination rule $K = \pi(B + F - \Delta)$ becomes a
cross-sector spectral invariant of the lattice-gauge structure:
$B$ is a Casimir-weighted matter trace on the $S^2$ base, $F$ is a
spectral zeta of the photon-sector modes in the infinite-shell
limit, and $\Delta$ is a Dirac-mode boundary count.  The additive
rule itself remains conjectural, just as in Paper~2; what changes
is the vocabulary available to talk about it.

What is sharp in this paper is the mathematics: the GeoVac graph
\emph{is} a discrete Hodge 1-complex, \emph{is} a Wilson lattice
gauge structure under the natural $U(1)$ edge-phase assignment, and
\emph{has} an exactly computable $S^2$ quotient that is Paper~2's
base graph.  What is conjectural is the physical interpretation of
the bundle components and the additive combination rule.  The
separation is deliberate and honest: the mathematical decomposition
needs no further justification, while the physical reading remains
Paper~2's conjecture.

The framework observation does not resolve the open problem of why
$K = \pi(B + F - \Delta)$ reproduces $\alpha^{-1}$.  It does locate
the three quantities in a single mathematical context, and it does
point to natural next questions: continuum-limit convergence of the
edge Laplacian, non-abelian generalizations on the Bargmann--Segal
$S^5$ lattice, higher-form analogs via the Breit rank-2 tensor, and
independent observables beyond $\alpha$.  These are the tractable
research directions that the observational synthesis makes
well-posed.

The broader takeaway is that the GeoVac project arrived at
lattice gauge structure without intending to: the requirements of a
sparse, $\pi$-free, quantum-number-labeled graph for computational
quantum chemistry happened to coincide, at the Hopf base of $S^3$,
with the incidence structure that lattice gauge theorists know as a
Wilson lattice with a $U(1)$ connection.  That the same graph
carries both content is either a coincidence of independent
constructions or a signal that the physical arena of Coulomb
bound-state quantum mechanics is, in its discrete form, exactly a
lattice gauge theory of matter coupled to $U(1)$ gauge fields on
the Hopf base.  The present paper does not adjudicate between these
two readings.  It documents that both are consistent with the
established mathematics, and it points to the follow-up computations
that could distinguish them.

%% ========================================================================
%% ACKNOWLEDGMENTS
%% ========================================================================

\begin{acknowledgments}
This paper was developed using an AI-augmented agentic workflow, with
implementation and documentation drafting performed collaboratively
with Anthropic's Claude.  The framework observation benefited from
the Phase~4B--4I alpha-sprint series, which systematically identified
the three spectral homes of the $(B, F, \Delta)$ ingredients over
approximately nine months of sustained computational decomposition.
\end{acknowledgments}

%% ========================================================================
%% REFERENCES
%% ========================================================================

\begin{thebibliography}{99}

\bibitem{Fock1935}
V.~Fock, Z.\ Phys.\ \textbf{98}, 145 (1935).

\bibitem{Wilson1974}
K.~G.~Wilson, Phys.\ Rev.\ D \textbf{10}, 2445 (1974).

\bibitem{Luscher1982}
M.~L\"uscher, Commun.\ Math.\ Phys.\ \textbf{85}, 39 (1982).

\bibitem{BergLuscher1981}
B.~Berg and M.~L\"uscher, Nucl.\ Phys.\ B \textbf{190}, 412 (1981).

\bibitem{Aharony2004}
O.~Aharony, J.~Marsano, S.~Minwalla, K.~Papadodimas, and M.~Van
Raamsdonk, Adv.\ Theor.\ Math.\ Phys.\ \textbf{8}, 603 (2004).

\bibitem{Witten1989}
E.~Witten, Commun.\ Math.\ Phys.\ \textbf{121}, 351 (1989).

\bibitem{Maldacena1998}
J.~M.~Maldacena, Adv.\ Theor.\ Math.\ Phys.\ \textbf{2}, 231 (1998).

\bibitem{Cunningham1909}
E.~Cunningham, Proc.\ London Math.\ Soc.\ \textbf{8}, 77 (1909).

\bibitem{Bateman1910}
H.~Bateman, Proc.\ London Math.\ Soc.\ \textbf{8}, 223 (1910).

\bibitem{EastwoodSinger1985}
M.~Eastwood and M.~Singer, Phys.\ Lett.\ A \textbf{107}, 73 (1985).

\bibitem{IkedaTaniguchi1978}
A.~Ikeda and Y.~Taniguchi, Osaka J.\ Math.\ \textbf{15}, 515 (1978).

\bibitem{Lim2020}
L.-H.~Lim, SIAM Review \textbf{62}, 685 (2020).

\bibitem{Schaub2020}
M.~T.~Schaub, A.~R.~Benson, P.~Horn, G.~Lippner, and A.~Jadbabaie,
SIAM Review \textbf{62}, 353 (2020).

\bibitem{Paper0}
J.~Loutey, ``Geometric packing and the origin of quantum numbers,''
GeoVac Paper~0 (2025).

\bibitem{Paper2}
J.~Loutey, ``The fine structure constant from Hopf bundle spectral
invariants,'' GeoVac Paper~2 (2025).

\bibitem{Paper7}
J.~Loutey, ``The dimensionless vacuum: Schr\"odinger recovery from the
unit $S^3$,'' GeoVac Paper~7 (2025).

\bibitem{Paper14}
J.~Loutey, ``Structurally sparse qubit Hamiltonians from natural
geometry,'' GeoVac Paper~14 (2026).

\bibitem{Paper18}
J.~Loutey, ``Spectral-geometric exchange constants: Projecting
discrete spectra onto continuous manifolds,'' GeoVac Paper~18 (2026).

\bibitem{Paper21}
J.~Loutey, ``The geometric vacuum: Quantum structure from compact
manifold topology,'' GeoVac Paper~21 (2026).

\bibitem{Paper22}
J.~Loutey, ``Potential-independent angular sparsity theorem,''
GeoVac Paper~22 (2026).

\bibitem{Paper24}
J.~Loutey, ``The Bargmann--Segal lattice: $\pi$-free discretization
of the 3D harmonic oscillator on $S^5$,'' GeoVac Paper~24 (2026).

\bibitem{Q1B}
GeoVac Track Q1-B, ``Edge Laplacian of the GeoVac $S^3$ graph at
$n_{\max} = 3$,'' internal memo, April 2026.

\end{thebibliography}

\end{document}
